\section{拓展}
\label{sec:future}

下面几项用于拓展当前初版架构：
\begin{itemize}[leftmargin=*,itemsep=2pt,topsep=2pt]
    \item \textbf{可训练的结构图：}使用端到端或单独训练的卷积模块，生成与输入空间尺寸对应的结构强度图。
    \item \textbf{主方向环扫描：}以每个等高线采样点为中心建立多层同心环，取代法线扫描。用局部主方向统一每个环的起点、返回点和汇集位置；先沿各环进行双向Mamba扫描，再按照从外到内的顺序把各环状态汇聚到中心。
    \item \textbf{层间稠密扫描：}在相邻等高线之间加入横向和纵向扫描，补充稀疏法线未读取的区域，再把稠密状态汇聚到附近的等高线采样点。等高线上也进行逐像素的稠密采样。
    \item \textbf{稠密层采样：}连续保留相邻层级，使层与层之间不留间隔。该方案取消法线扫描，只扫描稠密分布的等高线采样点，由相邻采样层形成连续覆盖。
\end{itemize}
